\documentclass{article}
\usepackage{graphicx}
\usepackage{amsmath}

\title{Lecture 4/13}
\author{Dylan Tang}
\date{April 2026}

\begin{document}

\maketitle

\section{Case Study: Vaping}
\begin{itemize}
    \item Conducted a randomized study that investigated vaping at UChicago Medicine
    \item $N_M = N_F = 112$
    \item $N_{M, \text{vape}} = 61$, $N_{F, \text{vape}} = 39$
\end{itemize}

 for proportions: \\
$H_0$: There is no difference in the proportion of males that vape versus the proportion of female that vape, i.e. $p_m = 0.5$. \\
$H_A$: The proportion of males that vape is greater than the proportion of females that vape. \\

\[
    C = \hat{p}_m \pm z^*_{1 - \frac{\alpha}{2}}\frac{\sigma_{pop}}{n}
\]

We can estimate $\sigma_{pop}$ with the standard error for proportions:
\[
     C = \hat{p}_M \pm z^*_{1 - \frac{\alpha}{2}} \sqrt{\frac{\hat{p}_M(1 - \hat{p}_M)}{n}}
\]
To simplify calculations, it is often to easier to invoke the upper bound bound by:
\[
    \sqrt{\frac{\hat{p}_M(1 - \hat{p}_M)}{n}} \leq \sqrt{\frac{0.5(1-0.5)}{n}}
\]
A conservative confidence interval can even be constructed knowing only the sample statistic and the range of the interval $[a,b]$.
\begin{itemize}
    \item The variance of the sample statistic must satisfy $Var(X_{stat}) \leq \frac{(b-a)^2}{4}$
    \item We can use this to construct a CI for sample proportions, if for some reason we don't know SE:
    \[
        C = \hat{p} \pm z^*_{1 - \frac{\alpha}{2}}\sqrt{\frac{(1-0)^2}{4}}
    \]
    \[
        \implies C = \hat{p} \pm \frac{1}{2}z^*_{1 - \frac{\alpha}{2}}
    \]
\end{itemize}

If $C$ captures 0.5, fail to reject $H_0$ at the $\alpha$ significance level, else reject.

Equivalently, we can conduct a one-sided test for proportions.

\[
    z = \frac{\hat{p}_{M} - 0.5}{\sqrt{\frac{0.5(1-0.5)}{n}}}
\]

If $z > 2$, reject at $\alpha = 0.5$ significance level.

\section{Case Study: Sensors}
\begin{itemize}
    \item There are four departments. Each department had 100 independent measurements.
    \item The averages for each department is the following:
    \begin{itemize}
        \item $\bar{x}_1 = 6.1$
        \item $\bar{x}_2 = 6.05$
        \item $\bar{x}_3 = 6.07$
        \item $\bar{x}_4 = 6.11$
    \end{itemize}
    \item $\hat{\mu}_{sensor} = 6.09$, $\sigma_{sensor} = 3.7$
    \item What is suspicious about this study?
    \begin{itemize}
        \item Standard deviation of sample means: $\frac{\sigma}{\sqrt{100}} = \frac{3.7}{10} = 0.37$. 
        \item But, the sample means observed above seem to be clustered much more closely than the standard deviation of sample means seem to suggest.
        \item To fully test, must run F-test with $H_0: \mu_1 = \mu_2 = \mu_3 = \mu_4$
    \end{itemize}
\end{itemize}

\section{The Linear Model}
The population is given by
\[
    Y = mX + b + \epsilon
\]
where $Y, X, \epsilon$ are R.Vs.

\subsection{Properties of the Population Model}
\begin{itemize}
    \item $E[\epsilon] = 0$
    \item $E[Y] = E[mX + b + \epsilon] = mE[X] + b$
    \item $Var(Y) = Var(mX + b + \epsilon) = \underbrace{m^2Var(X)}_{\text{signal}} + \underbrace{Var(\epsilon)}_{\text{noise}}$
    \item \textbf{Signal-Noise Ratio} = $\frac{m^2Var(X)}{Var(\epsilon)}$
\end{itemize}
\end{document}
